\usepackage{fontspec}
\usepackage[brazil]{babel}
\usepackage{microtype}
\usepackage[a4paper,margin=2.6cm,headheight=15pt]{geometry}
\usepackage{graphicx}
\usepackage{booktabs}
\usepackage{longtable}
\usepackage{amsmath}
\usepackage[dvipsnames]{xcolor}
\usepackage{listings}
\usepackage[most]{tcolorbox}
\usepackage{fancyhdr}
\usepackage{titlesec}
\usepackage{natbib}
\usepackage{makeidx}
\usepackage{hyperref}

% Fontes distribuídas normalmente com o TeX Live.
\setmainfont{TeX Gyre Pagella}
\setsansfont{TeX Gyre Heros}
\setmonofont{TeX Gyre Cursor}
\setlength{\emergencystretch}{1em}

\definecolor{AzulLivro}{HTML}{183B56}
\definecolor{DouradoLivro}{HTML}{D19A32}
\definecolor{FundoCodigo}{HTML}{F5F7F9}
\definecolor{VerdeCodigo}{HTML}{356859}

\hypersetup{
  colorlinks=true,
  linkcolor=AzulLivro,
  citecolor=VerdeCodigo,
  urlcolor=MidnightBlue,
  pdftitle={LaTeX do Zero ao Primeiro Livro}
}

\pagestyle{fancy}
\fancyhf{}
\fancyhead[LE]{\small\thepage\quad\nouppercase{\leftmark}}
\fancyhead[RO]{\small\nouppercase{\rightmark}\quad\thepage}
\renewcommand{\headrulewidth}{0.4pt}

\titleformat{\chapter}[display]
  {\sffamily\bfseries\color{AzulLivro}}
  {\filleft\Large\chaptername\ \thechapter}
  {1ex}
  {\titlerule\vspace{1ex}\Huge}

\lstset{
  basicstyle=\ttfamily\small,
  backgroundcolor=\color{FundoCodigo},
  keywordstyle=\color{AzulLivro}\bfseries,
  commentstyle=\color{VerdeCodigo}\itshape,
  stringstyle=\color{BrickRed},
  numbers=left,
  numberstyle=\tiny\color{gray},
  frame=single,
  rulecolor=\color{gray!30},
  breaklines=true,
  showstringspaces=false,
  tabsize=2
}

\lstdefinelanguage{LaTeXBook}{
  language=[LaTeX]TeX,
  morekeywords={documentclass,usepackage,begin,end,chapter,section,
    includegraphics,label,ref,cite,frontmatter,mainmatter,backmatter,include}
}

\newtcolorbox{pratica}[1][]{
  colback=blue!3,colframe=AzulLivro,fonttitle=\bfseries,
  title=Pratique agora,#1
}
\newtcolorbox{atencao}[1][]{
  colback=yellow!8,colframe=DouradoLivro!80!black,fonttitle=\bfseries,
  title=Atenção,#1
}
\newtcolorbox{resultado}[1][]{
  colback=green!4,colframe=VerdeCodigo,fonttitle=\bfseries,
  title=Resultado esperado,#1
}

\newcommand{\comando}[1]{\texttt{\textbackslash#1}}
\newcommand{\arquivo}[1]{\texttt{#1}}
\makeindex
